\section*{Noncentral Holonomy Versus Brauer-Type Projective Twists}

Central/projective obstructions and noncentral nonabelian holonomy are different
phenomena.  A Brauer-type or projective torsion diagnostic is looking for a
central scalar residual, for example a commutator equal to a root of unity times
the identity.  A permutation-valued residual can instead be noncentral inside
the generated permutation group.  Such a residual may obstruct naive descent or
synchronization, but it is not by itself a scalar Brauer/projective class.

\paragraph{Example.}
In \(S_3\), let \(p=(12)\) and \(q=(23)\).  The commutator
\([p,q]=p q p^{-1}q^{-1}\) is a nontrivial 3-cycle.  Since the center of
\(S_3\) is trivial, this 3-cycle is not central in \(S_3\).  Therefore this
obstruction is a noncentral permutation holonomy example, not a scalar
Brauer/projective class.

Pure permutation synchronization, as in C2M3-style alignment, is an untwisted
or nonabelian gauge-synchronization problem.  TwistedMerge++ extends this by
testing whether the residual descends to a central scalar torsion class after
the structure group has been enlarged to signed, phase, monomial, block, or
projective gauges.

Failure to find finite-index scalar residuals in permutation alignments does
not mean all obstruction theory is trivial.  It means that the detected
residuals are noncentral in the chosen structure group, or that a larger
structure group must be tested before making Brauer/projective claims.
